\documentclass[12pt]{article}
\usepackage[T1]{fontenc}
\usepackage{lmodern}
\usepackage[margin=1in]{geometry}
\usepackage{natbib}
\usepackage{amsmath}
\usepackage{booktabs}
\usepackage{threeparttable}
\usepackage{setspace}
\onehalfspacing
\title{Contracts and Reminders at the Gym: A Field Experiment on Commitment Devices for Exercise\thanks{We thank seminar participants at Example University and the staff of the campus recreation center for their help. The experiment was approved by the Institutional Review Board of Example University (protocol 2024-0391) and registered before enrollment in the American Economic Association's registry for randomized controlled trials (AEARCTR-0098317). Funding came from the Example University Wellness Research Fund, which had no role in the design, the analysis or the decision to publish. The authors have no conflicts of interest to declare.}}
\author{Alex Example\thanks{Department of Economics, Example University, alex@example.edu} \and Riley Sample\thanks{School of Public Health, Example University}}
\date{October 2026}
\begin{document}
\maketitle
\begin{abstract}
Commitment contracts and reminders are cheap tools against procrastination, but they are rarely tested together. We randomized 118 university departments, with 2,412 employees who held a campus gym membership, to a control group, text-message reminders, the offer of a self-funded commitment contract, or both. A third of the employees offered a contract signed one. The offer raised gym visits by 37 percent during the eight weeks in which the contracts ran. Reminders added little on their own but made the offer more effective. Effects shrank once the contracts expired. The effect is concentrated among employees who used the gym less than once a week before the study. We discuss what the results imply for workplace wellness programs.
\end{abstract}

\noindent JEL codes: C93, D91, I12

\noindent Keywords: commitment devices, reminders, exercise, present bias, field experiment

\section{Introduction}
\label{sec:intro}
Many people who pay for a gym membership use it far less than they had planned \citep{dellavigna2006}. Models of present bias explain this pattern: a person who gives extra weight to immediate costs keeps postponing an activity whose costs come now and whose benefits come later, and she may not foresee that she will keep postponing it \citep{laibson1997,odonoghue1999}. Paying people to exercise raises attendance while the payments last, but the effects often fade once the payments stop \citep{gneezy2012}.

Two cheap tools address the problem from different sides. A commitment contract raises the cost of skipping the gym for a person who expects to be tempted later \citep{bryan2010}. A reminder makes a planned activity salient at the moment of choice, which helps people who are forgetful rather than impatient \citep{karlan2016}. Evidence on each tool is growing, but the two are rarely tested together, even though the case for combining them is intuitive: a contract gives people a reason to go, and a reminder tells them when.

This paper reports a field experiment with 2,412 employees of a large public university who held a membership at the campus recreation center. We randomized their departments to four arms in a two-by-two design: a control group, text-message reminders, the offer of a self-funded commitment contract, and the offer combined with reminders. A contract required at least 12 visits in eight weeks. Participants chose their own stake, which was refunded if they met the target and donated to a food bank if they did not. Attendance comes from the recreation center's entry records for the 12 weeks before and the 20 weeks after the start of the intervention.

Offering a commitment contract caused a large increase in attendance while the contracts were in force. About a third of the employees offered a contract signed one, and the offer raised visits by 0.31 per week. Reminders did little on their own but made the offer more effective. After the contracts expired, the effect of the offer on visits fell to about a third of its size.

The paper makes three contributions. First, it tests whether reminders and commitment contracts are complements, in a population that has already taken the first step of buying a membership. Second, it randomizes departments rather than individuals, which limits spillovers among colleagues who exercise together. Third, it relates the take-up of contracts to stated self-control problems and compares the cost of the contracts with their effects.

Section~\ref{sec:literature} reviews related work. Section~\ref{sec:design} describes the design and Section~\ref{sec:data} the data. Section~\ref{sec:results} reports the results and Section~\ref{sec:robustness} robustness checks. Section~\ref{sec:discussion} discusses mechanisms and costs, and Section~\ref{sec:conclusion} concludes.

\section{Related literature}
\label{sec:literature}
Our study builds on three strands of work: incentives for exercise and the habits they may create, commitment devices, and reminders and plans.

\subsection{Incentives and habits}
A first strand pays people to exercise and asks whether a habit survives once the payments stop. \citet{charness2009} find lasting increases in attendance after a month of paid gym visits, driven by people who had not been regular gym users before. \citet{acland2015} replicate this design with students and staff of a large public university and find that the habit formed during the incentive month largely decayed over the winter break. \citet{royer2015} pay employees of a Fortune 500 company for visits to the onsite gym and find strong responses while the payments last but little persistence, unless the incentive is followed by the offer of a commitment contract. Among new members of a private gym, \citet{carrera2018incentives} find that modest threshold incentives had small effects on attendance, and that the effects on average visits were largest among members who already exercised regularly.

\subsection{Commitment devices}
A second strand studies commitment devices directly. Theory predicts a demand for commitment from people who are aware of their present bias \citep{odonoghue1999}, and \citet{bryan2010} review commitment devices in saving, health and other domains. \citet{augenblick2015} find considerable present bias in the allocation of effort but very little in choices over money, and 59 percent of their subjects chose a commitment device when it was offered at a price of zero. In a data-entry firm in India, workers chose contracts that penalized them for low output \citep{kaur2015}. Deposit contracts, in which participants put their own money at risk, helped obese adults lose weight in a randomized trial \citep{volpp2008}.

Commitment contracts for exercise have been studied in workplaces and online. In \citet{royer2015}, 22 percent of the employees offered a self-funded commitment contract signed one. Those who committed placed 58 dollars at stake on average, and 63 percent of them met their target. \citet{bhattacharya2015} randomize the default length of exercise contracts on the stickK.com platform. Longer defaults led users to choose longer contracts but also made them less likely to sign a contract at all. Only 22.7 percent of the signed contracts put money at stake.

Commitment does not have to involve money. \citet{milkman2014} let students, faculty and staff of a university listen to tempting audiobooks only at the gym. Attendance rose until the Thanksgiving break, and the program helped the busiest participants most. When offered the chance to pay for a month of gym-only access to a tempting audiobook, participants in that study were willing to pay \$9.61 on average.

Recent work questions whether the take-up of contracts reveals awareness of self-control problems. In an experiment with 1,248 members of a fitness center, \citet{carrera2022} find high take-up of contracts that pay only if attendance reaches a threshold (64, 49 and 32 percent for thresholds of 8, 12 and 16 visits in four weeks), but an information treatment that made members more aware of their overoptimism reduced demand for contracts. They estimate an average present-bias parameter of 0.84 and a perceived parameter of 0.55, so that members are more pessimistic about their future self-control than their behavior warrants.

\subsection{Reminders and plans}
A third strand studies reminders and plans. Reminders help clients of commitment savings accounts reach their savings goals \citep{karlan2016}, and prompts to write down the date and time of a planned flu shot raise vaccination rates \citep{milkman2011}. For repeated behavior the evidence is weaker. \citet{carrera2018plans} asked members of a private gym to plan their visits for the next two weeks and find a precisely estimated null effect on attendance; total planned visits were three times actual visits (7.04 against 2.33). In the megastudy of \citet{milkman2021}, which tested 54 four-week programs with 61,293 members of a gym chain, a basic condition combining planning prompts, reminders and small rewards raised weekly visits by 0.14, or 9 percent. The authors report unadjusted p-values for their 53 comparisons together with an estimate of the false discovery rate.

We add a direct test of whether reminders make commitment contracts more effective, with contracts that participants fund themselves and with randomization across departments rather than individuals.

\section{Experimental design}
\label{sec:design}

\subsection{Setting and recruitment}
The study took place at a large public university in the Midwest. Employees can join the campus recreation center for 28 dollars a month, deducted from their pay. The center records each entry with an ID-card swipe. Members do not swipe out, so we count visits but cannot measure how long they last. \citet{acland2015}, by contrast, recorded both entry and exit swipes and could therefore measure how long each visit lasted. \citet{carrera2018incentives} introduced a swipe-out requirement partway through their study and credited a visit toward the incentive only after at least ten minutes in the gym.

In January 2025 we invited by email all 7,940 employees with an active membership to take part in a study of exercise habits. Of these, 2,537 gave informed consent and completed a baseline survey, for which they received a 10-dollar gift card. We excluded 125 employees who left the university or let their membership lapse before the start date, leaving 2,412 participants in 118 departments.

Enrollment and the baseline survey ran from 6 to 24 January 2025. We randomized departments on 27 January, and the intervention began on Monday 3 February 2025. The eight intervention weeks ended on 30 March, and the twelve follow-up weeks on 22 June. Weeks run from Monday to Sunday. A visit is a day with at least one entry swipe, so a second swipe on the same day does not count as a second visit, and the contract target of 12 visits refers to 12 different days.

\subsection{Randomization}
We randomized departments rather than individuals. Colleagues often exercise together and talk about workplace programs, and randomizing whole departments limits spillovers between arms. Randomization was stratified by department size (above or below 20 participants) and by the department's mean baseline attendance (above or below the median), which gives four strata. As Table~\ref{tab:balance} shows, 30, 29, 30 and 29 departments were assigned to the control, reminder, commitment and combined arms. All participants in a department received the same treatment.

\subsection{Treatments}
All arms received a one-page leaflet on the health benefits of regular exercise.

\emph{Reminders.} In the week before the intervention began, participants in the two reminder arms chose up to three weekdays for exercise on a short online form. During the eight intervention weeks they received a text message at 7 a.m. on each chosen day, for example: ``Today is one of your gym days. The recreation center opens at 6 a.m.''

\emph{Commitment contracts.} Participants in the two commitment arms received a contract offer by email in the first intervention week. The contract required at least 12 visits during the eight intervention weeks. Participants chose a stake between 20 and 100 dollars and paid it by card. The stake was refunded after week 8 if the participant met the target and was donated to a local food bank if not. Declining the offer had no consequences, and participants who signed could not change their stake later. The offer email explained the rules in four sentences and linked to a one-page form on which participants entered their stake. Participants who had not answered after four days received one follow-up email, and the offer closed at the end of the first intervention week. The staff of the recreation center did not know which participants held contracts.

\emph{Control.} The control group received the leaflet and nothing else.

No participant was paid for visiting the gym.

\subsection{Outcomes and registration}
We registered the design, the outcomes and the main specification before enrollment began. Like \citet{milkman2021}, who preregistered their megastudy as a single study before data collection, we registered a single primary outcome. The pre-registered primary outcome is the share of the eight contract weeks in which a participant visited the gym at least once; the number of visits per week is a secondary outcome. We measure persistence with the same two outcomes over weeks 9 to 20.

With 118 departments and about 600 participants per arm, a standard deviation of 0.9 visits per week and an intraclass correlation of 0.02 within departments, the design had 80 percent power at the 5 percent level for an MDE of 0.17 visits per week in any pairwise comparison of arms.

\subsection{Estimation}
We estimate intention-to-treat (ITT) effects with the regression
\begin{equation}
y_{it} = \alpha + \beta_1 C_i + \beta_2 R_i + \beta_3 C_i R_i + \gamma_t + \delta_{s(i)} + \varepsilon_{it},
\label{eq:main}
\end{equation}
where $y_{it}$ is an outcome of participant $i$ in week $t$, $C_i$ and $R_i$ indicate that the participant's department was assigned to the contract offer and to reminders, $\gamma_t$ are week fixed effects and $\delta_{s(i)}$ are fixed effects for the randomization stratum of the participant's department. The coefficient $\beta_1$ is the effect of the contract offer without reminders, $\beta_2$ the effect of reminders without the offer, and $\beta_3$ the change in the effect of the offer when reminders are added. Standard errors are clustered by participant because each participant contributes several weekly observations.

\section{Data}
\label{sec:data}
Table~\ref{tab:balance} reports baseline characteristics by arm. Sixty-one percent of participants are women, and their average age is 41. In the 12 weeks before the study, participants made 0.86 visits per week on average and visited the gym in 39 percent of weeks. About two thirds said at baseline that they exercised less often than their own target. The arms are similar on every characteristic, and none of the joint tests in the last column of Table~\ref{tab:balance} rejects equality at conventional levels.

Attendance comes from the recreation center's entry records for every randomized participant, so no outcome data are missing. The panel covers the 12 baseline weeks and the 20 weeks from the start of the intervention.

Of the 1,203 employees offered a contract, 409 signed one (34.0 percent): 198 of 603 in the commitment arm (32.8 percent) and 211 of 600 in the combined arm (35.2 percent). The mean stake was 46 dollars and the median 40 dollars. Of the 409 signers, 241 (58.9 percent) met the target and got their stake back.

\begin{table}[t]
\centering
\small
\begin{threeparttable}
\caption{Baseline characteristics by arm}
\label{tab:balance}
\begin{tabular}{@{}p{5.2cm}ccccc@{}}
\toprule
 & Control & Reminders & Commitment & Commitment & p-value \\
 & & & & and reminders & \\
\midrule
Female & 0.612 & 0.601 & 0.619 & 0.605 & 0.93 \\
Age (years) & 41.3 & 40.8 & 41.0 & 41.6 & 0.81 \\
Position: faculty & 0.284 & 0.291 & 0.277 & 0.288 & 0.89 \\
Position: staff & 0.598 & 0.603 & 0.609 & 0.597 & 0.97 \\
Position: other & 0.118 & 0.124 & 0.114 & 0.115 & 0.95 \\
Visits per week, 12 weeks before the study & 0.85 & 0.87 & 0.86 & 0.87 & 0.84 \\
Share of weeks with a visit, 12 weeks before the study & 0.391 & 0.396 & 0.388 & 0.399 & 0.92 \\
Exercises less than own target & 0.684 & 0.679 & 0.690 & 0.677 & 0.95 \\
Minutes from office to the gym & 9.6 & 9.9 & 9.4 & 9.7 & 0.77 \\
\midrule
Participants & 598 & 611 & 603 & 600 & \\
Departments & 30 & 29 & 30 & 29 & \\
\bottomrule
\end{tabular}
\begin{tablenotes}
\small
\item[] Notes: Means of baseline characteristics by arm. Shares are fractions of participants. The last column reports the p-value of a joint test that the four arm means are equal, from a regression on arm indicators with stratum fixed effects and standard errors clustered by department.
\end{tablenotes}
\end{threeparttable}
\end{table}


\section{Results}
\label{sec:results}

\subsection{Effects during the contract period}
Table~\ref{tab:main} reports ITT estimates of equation~\eqref{eq:main}. The estimates use all 2,416 randomized employees and their weekly attendance in the eight intervention weeks (columns 1 and 2) and in the twelve weeks that followed (columns 3 and 4).

We focus on column (2), which reports the primary outcome, the number of visits per week. Offering a contract without reminders raised attendance by 0.31 visits per week (95 percent confidence interval 0.17 to 0.45), from 0.84 to 1.15 visits per week, an increase of 37 percent. The point estimate for reminders alone is 0.064 visits per week (standard error 0.052). This estimate is small and not statistically significant at conventional levels. The interaction term is positive: reminders raise the effect of the contract offer by 0.152 visits per week. Employees offered both a contract and reminders made 0.51 more visits per week than the control group, an increase of 57 percent.

Because the regression includes the interaction, the coefficient on the contract offer measures its effect in the absence of reminders, and the coefficient on reminders measures their effect in the absence of an offer. The effect of the offer when reminders are also sent is the sum of the first and third coefficients, 0.462 visits per week.

Column (1) shows the same pattern for the share of weeks with at least one visit. The contract offer raised this share by 6.5 percentage points, or 17 percent of the control mean of 38.0 percent. The estimates for reminders alone and for the interaction are small and imprecise for this outcome.

\begin{table}[t]
\centering
\small
\begin{threeparttable}
\caption{Effects of contract offers and reminders on gym attendance}
\label{tab:main}
\begin{tabular}{@{}p{5.7cm}cccc@{}}
\toprule
 & (1) & (2) & (3) & (4) \\
 & \multicolumn{2}{c}{Weeks 1--8} & \multicolumn{2}{c}{Weeks 9--20} \\
 & Any visit & Visits per week & Any visit & Visits per week \\
\midrule
Commitment offer & $0.065^{***}$ & $0.310^{***}$ & $0.019$ & $0.112^{**}$ \\
 & (0.017) & (0.071) & (0.013) & (0.059) \\
Reminders & $0.012$ & $0.046$ & $-0.003$ & $0.009$ \\
 & (0.015) & (0.052) & (0.012) & (0.044) \\
Commitment offer × Reminders & $0.024$ & $0.152^{**}$ & $0.011$ & $0.038$ \\
 & (0.024) & (0.076) & (0.019) & (0.064) \\
\midrule
Control mean & 0.380 & 0.84 & 0.335 & 0.71 \\
Observations (participant-weeks) & 19,296 & 19,296 & 28,944 & 28,944 \\
Participants & 2,412 & 2,412 & 2,412 & 2,412 \\
\bottomrule
\end{tabular}
\begin{tablenotes}
\small
\item[] Notes: Intention-to-treat estimates of equation~\eqref{eq:main} from linear regressions of weekly outcomes on the arm indicators, week fixed effects and randomization-stratum fixed effects. Any visit is an indicator for at least one visit in the week, so its mean is the share of weeks with a visit. Weeks 1--8 are the contract period and weeks 9--20 the twelve weeks after it. Control mean is the mean of the outcome in the control arm over the same weeks. Standard errors clustered by participant in parentheses. * $p<0.10$, ** $p<0.05$, *** $p<0.01$.
\end{tablenotes}
\end{threeparttable}
\end{table}


\subsection{After the contracts expired}
Columns (3) and (4) cover weeks 9 to 20. The effect of the contract offer on visits falls to 0.112 per week, about a third of its size during the contract period, and remains significant at the 5 percent level. The share of weeks with a visit is 1.9 percentage points higher in the commitment arm, with a 95 percent interval of 0.4 to 3.4 points. The reminder estimates are close to zero after week 8, and so are the interaction terms.

\subsection{Signers and decliners}
How much of the effect comes from the contracts themselves? During the eight contract weeks, signers made 1.94 visits per week and decliners 0.89. Signing a contract therefore more than doubled gym attendance. Signers had also been more frequent gym users before the study, with 1.18 visits per week at baseline against 0.70 for decliners. Random assignment offers a second way to scale the effect: dividing the effect of the offer in the arm without reminders (0.310) by the take-up rate in that arm (32.8 percent) gives about 0.95 visits per week for those who signed.

\subsection{Who signs a contract}
Take-up was higher among employees who said at baseline that they exercised less often than their own target: 37.1 percent against 27.3 percent. This pattern is consistent with \citet{royer2015}, who find that the demand for commitment comes mainly from employees who fall short of their own exercise targets. Women were more likely than men to sign (36.3 against 30.4 percent). These take-up correlates were not pre-specified and are descriptive.

\subsection{Heterogeneity}
Table~\ref{tab:het} reports the effect of the contract offer in five pairs of subgroups, pooling the arms with and without reminders. The pooled effect is 0.386 visits per week, the average of the effects of the offer with and without reminders in Table~\ref{tab:main}. The commitment offer works mainly for employees who rarely used the gym before the study: among employees with fewer than one weekly visit at baseline, it raised visits by 0.508 per week, against 0.187 among more frequent users (p = 0.040 for the difference). The other four splits, by gender, age, position and travel time to the gym, show no significant differences.

\begin{table}[t]
\centering
\small
\begin{threeparttable}
\caption{Effect of the contract offer on visits per week by subgroup, weeks 1--8}
\label{tab:het}
\begin{tabular}{@{}p{6.4cm}ccc@{}}
\toprule
 & (1) & (2) & (3) \\
Split (group A versus group B) & Group A & Group B & p-value, A = B \\
\midrule
Women versus men & $0.431^{***}$ & $0.316^{***}$ & 0.36 \\
 & (0.078) & (0.098) & \\
Under 40 versus 40 and over & $0.352^{***}$ & $0.415^{***}$ & 0.61 \\
 & (0.092) & (0.083) & \\
Faculty versus staff and other & $0.288^{**}$ & $0.425^{***}$ & 0.33 \\
 & (0.121) & (0.071) & \\
Fewer than one baseline visit per week versus one or more & $0.508^{***}$ & $0.187$ & 0.040 \\
 & (0.088) & (0.129) & \\
Under 10 minutes to the gym versus 10 minutes or more & $0.398^{***}$ & $0.371^{***}$ & 0.83 \\
 & (0.081) & (0.094) & \\
\midrule
All participants & $0.386^{***}$ & & \\
 & (0.061) & & \\
\bottomrule
\end{tabular}
\begin{tablenotes}
\small
\item[] Notes: Each row reports the effect of the contract offer, pooling the arms with and without reminders, estimated separately in the two groups of a split, with the specification of Table~\ref{tab:main}. Column (3) reports the p-value of a test that the two effects are equal; p-values are not adjusted for multiple comparisons. Standard errors clustered by participant in parentheses. * $p<0.10$, ** $p<0.05$, *** $p<0.01$.
\end{tablenotes}
\end{threeparttable}
\end{table}


\section{Robustness}
\label{sec:robustness}
Table~\ref{tab:robust} shows that the main estimate is stable across specifications. Adding baseline controls for prior attendance, age, gender and position lowers the estimate slightly, to 0.297 visits per week, and makes it more precise (column 2). Because visits are counts, column (3) estimates a Poisson model; the Poisson coefficient of 0.312 implies an increase of 37 percent, close to the linear estimate. Dropping the first intervention week, in which participants were still signing contracts, gives 0.301 (column 4). Dropping the nine departments with fewer than five participants, which together have 31 participants, gives 0.318 (column 5).

Two further checks are not shown in the table. The campus was closed for two days in week 6, and dropping that week leaves the estimate at 0.31. Excluding the 37 participants who moved to another department during the study barely changes the estimate (0.309).

\begin{table}[t]
\centering
\small
\begin{threeparttable}
\caption{Robustness of the effect of the contract offer on visits per week, weeks 1--8}
\label{tab:robust}
\begin{tabular}{@{}p{4.8cm}ccccc@{}}
\toprule
 & (1) & (2) & (3) & (4) & (5) \\
 & Main & Baseline & Poisson & Without & Without small \\
 & & controls & & week 1 & departments \\
\midrule
Commitment offer & $0.310^{***}$ & $0.297^{***}$ & $0.312^{***}$ & $0.301^{***}$ & $0.318^{***}$ \\
 & (0.071) & (0.058) & (0.078) & (0.074) & (0.075) \\
\midrule
Observations (participant-weeks) & 19,296 & 19,296 & 19,296 & 16,884 & 19,048 \\
Participants & 2,412 & 2,412 & 2,412 & 2,412 & 2,381 \\
\bottomrule
\end{tabular}
\begin{tablenotes}
\small
\item[] Notes: The dependent variable is the number of visits in a week. All columns include the reminder indicator, its interaction with the contract offer, week fixed effects and stratum fixed effects; only the coefficient on the contract offer is shown. Column (2) adds baseline visits per week, age, gender and position. Column (3) reports a Poisson coefficient in log points. Column (5) drops the nine departments with fewer than five participants. Standard errors clustered by participant in parentheses. * $p<0.10$, ** $p<0.05$, *** $p<0.01$.
\end{tablenotes}
\end{threeparttable}
\end{table}


\section{Discussion}
\label{sec:discussion}

\subsection{Why do reminders help only with a contract?}
Reminders on their own do not change behavior; they work only as a complement to the contract. A reminder tells a participant that today is a gym day, but without a contract skipping the visit costs nothing. With a contract, the reminder arrives while the stake is at risk, which may be why the combination works better than either tool alone. The pattern is also consistent with limited attention to a distant target: a contract holder who forgets the target until the last week may no longer be able to reach it.

\subsection{What does take-up reveal?}
Take-up need not reveal sophistication about present bias. \citet{carrera2018plans} find that between 27 and 34 percent of gym members accept a contract that pays them only if they visit the gym less often, which no member trying to exercise more should want. In our setting, a third of those offered a contract signed it, and 41 percent of the signers missed the target and lost their stake. Some of these signers may have overestimated their future self-control, as the partial naivet\'e documented by \citet{acland2015} would predict.

\subsection{Costs}
Administering the contracts, including staff time and card fees, cost 9.80 dollars per employee offered a contract. Since the offer raised visits by 0.31 per week over eight weeks, or 2.48 visits per employee offered a contract, each additional visit cost about 2.95 dollars. The forfeited stakes went to the food bank and are a transfer rather than a cost.

\subsection{Scope and limitations}
Our participants had already bought a membership, so the estimates speak to people who have taken a first step toward exercising more. Our data also have limits. We observe visits but not what people do at the gym, the sample comes from a single employer, and the follow-up ends twelve weeks after the contracts expired.

\section{Conclusion}
\label{sec:conclusion}
Offering employees a self-funded commitment contract raised gym attendance substantially while the contracts ran, and text-message reminders made the offer more effective. After the contracts expired, much of the effect faded. A third of the employees offered a contract signed one, and many of those who signed lost their stake.

For employers, the contracts are cheap to run and easy to explain. Their main weakness is that they ask employees to put their own money at risk, and four in ten signers lost theirs. A contract that returned part of a forfeited stake, or that let participants choose a lower target, might keep most of the effect at a lower private cost.

These findings suggest that commitment contracts and reminders work best together and that their effects last little longer than the contracts themselves. Future work could test whether renewable contracts sustain attendance, and whether contracts reach people who do not yet hold a gym membership. Our results suggest that health insurers and employers everywhere should offer commitment contracts with text-message reminders to all adults who want to exercise more.

\section*{Data availability}
The de-identified data and the code will be deposited at openICPSR when the paper is accepted for publication. The recreation center shares attendance records under a data use agreement that does not allow the public release of identifiable data.


\bibliographystyle{aer}
\bibliography{refs}
\end{document}
